\section{Problem 3.}
The two distinct lines with equations $ax + y = c$ and $x + y = c$
intersect at the point $(b - 2, 4)$.  Solve for $b + c$.

\begin{solution}
    Okay, this problem seems much easier than the previous ones. From the wording alone, it seems
    to hint that $a$ should be arbitriary with the sole exception of the $a \neq 1$ to satisfy the 
    distinct line requirement. Let us represent the system of equations with matrices and into a linear
    combination of vectors.

    \begin{align*}
        ax + y &= c \\
        x + y &= c \\
 \Longrightarrow
        \begin{bmatrix}
            a & 1 \\
            1 & 1 
        \end{bmatrix}
        \begin{bmatrix}
            x \\
            y
        \end{bmatrix}
        &=
        \begin{bmatrix}
            c \\
            c
        \end{bmatrix} \\
        \begin{bmatrix}
            a \\
            1
        \end{bmatrix} x
        + 
        \begin{bmatrix}
            1 \\
            1
        \end{bmatrix} y
        &= \begin{bmatrix}
            c \\
            c
        \end{bmatrix}
    \end{align*}

    Since the vectors are both linearly independent or in other
    words the null space is the zero vector, we are guarenteed only
    one solution when $x = 0, \ y = c$. We can show this with Column-Row factorization of the augmented matrix.

    \begin{align*}
        & 
        \begin{bmatrix}
            a & 1 & c \\
            1 & 1 & c
        \end{bmatrix}
        =
        \begin{bmatrix}
            a & 1 \\
            1 & 1 
        \end{bmatrix}
        \begin{bmatrix}
            1 & 0 & 0 \\
            0 & 1 & c
        \end{bmatrix}
    \end{align*}
    
    The factorization shows us also the row reduced echelon form which gets us an 
    intersection point at $(0, c)$. We simply just need to match this with our given
    intersection point with $(b-2, 4)$. 

    \begin{align*}
        & (0, c) = (b-2, 4) \\
        & 0 = b-2 \\
        & b = 2 \\
        & c = 4 \\
        \Longrightarrow & \boxed{b + c = 6}
    \end{align*}

\end{solution}

